\begin{resumo}
Avalia-se, por simulação, se a atribuição conjunta de pedidos em lotes de 60 segundos pelo algoritmo Húngaro reduz a distância percorrida sem passageiro até o embarque (km vazio), preservando o atendimento. Comparam-se o despacho imediato pelo veículo elegível mais próximo, o despacho sequencial em lote e a atribuição conjunta, sob os mesmos pedidos, veículos, jornadas e custos de distância. A avaliação principal utiliza dados históricos de uma empresa de transporte por aplicativo na Grande Vitória (ES); cenários com demanda sintética sobre a geografia real oferecem evidência exploratória separada. Com os parâmetros de março congelados antes das execuções, a atribuição conjunta reduziu o km vazio simulado em 5,96\% em setembro e 5,04\% em outubro de 2026 frente ao despacho imediato. Contudo, em outubro atendeu 85 pedidos a menos e não cumpriu o critério de atendimento. Frente ao despacho sequencial em lote, as reduções foram de 2,74\% e 1,40\%. Em teste posterior de 12 dias, com nova calibração, a diferença frente ao despacho imediato foi de −2,79\%, com intervalo que inclui zero e 18 atendimentos a menos; frente ao lote sequencial, houve 0,87\% mais km e 29 atendimentos a mais. Uma comparação posterior com um guloso global não sustenta superioridade geral da atribuição conjunta. O simulador tampouco reproduziu o atendimento observado na operação. Assim, os resultados caracterizam diferenças entre políticas no modelo e não demonstram economia operacional.
\end{resumo}

\begin{abstract}
Simulation is used to assess whether jointly assigning requests in 60-second batches with the Hungarian algorithm reduces the distance driven without a passenger to pickup (empty kilometers) while preserving service. Immediate dispatch to the nearest eligible vehicle, sequential batch dispatch, and joint assignment are compared using the same requests, vehicles, shifts, and distance costs. The main evaluation uses historical data from a ride-hailing company in Greater Vitória, Brazil; synthetic demand over the real geography provides separate exploratory evidence. With March parameters frozen before execution, joint assignment reduced simulated empty kilometers by 5.96\% in September and 5.04\% in October 2026 relative to immediate dispatch. However, it served 85 fewer requests in October and failed the service criterion. Relative to sequential batch dispatch, reductions were 2.74\% and 1.40\%. In a later 12-day test with new calibration, the difference relative to immediate dispatch was −2.79\%, with an interval including zero and 18 fewer requests served; relative to sequential batch dispatch, joint assignment incurred 0.87\% more empty kilometers while serving 29 more requests. A later comparison with a global greedy rule does not establish general superiority of joint assignment. The simulator also failed to reproduce service observed in the operation. The results therefore describe differences among policies within the model and do not demonstrate operational savings.
\end{abstract}
